%----------------------------------------------------------------------------------------
% LIMIET VAN EEN RIJ. Alle opgaven zijn eigen werk.
%
% Opbouw: Start, leerstof (kern, uitgewerkt voorbeeld, korte oefening), en daarna op een
% nieuwe pagina de Opdrachten (door elkaar, stijgende moeilijkheid), Koppel en Test jezelf.
% Labels van dit hoofdstuk beginnen met h4-. \minigrafiek staat in hoofdstuk 1.
%----------------------------------------------------------------------------------------

\mychapter{4}{Limiet van een rij}

\lkr{Leerplan B+S'': LPD 45 (formele definitie, met kwantoren, toegepast op een concreet voorbeeld) en LPD 46 (limieten bepalen, convergentie, rekenregels, insluitstelling). De meetkundige reeks is een suggestie bij LPD 44 en staat hier omdat ze de limiet nodig heeft. Voorkennis: limieten van functies op oneindig. Rode draad: je kunt de limiet op drie manieren vinden, uit de grafiek, met rekenregels en uit een vergelijking, maar elke manier heeft haar voorwaarden.}

\begin{oefeningen}[Start]
\begin{oefening}
\blokkop{Opwarmer}
\begin{enumerate}
  \item Bereken $\displaystyle\lim_{x\to+\infty} \frac{1}{x}$, $\displaystyle\lim_{x\to+\infty} \frac{2x+1}{x-3}$ en $\displaystyle\lim_{x\to+\infty} (1 - x^2)$.
  \item Schrijf $|x - 5| < 0{,}1$ als een open interval.
  \item Voor welke $n > 0$ is $\dfrac{1}{n} < 0{,}001$?
  \item Voor welke $n$ is $0{,}5^{\,n} < 0{,}01$?
\end{enumerate}
\begin{solution}
\textbf{a.}~0; 2; $-\infty$ \quad \textbf{b.}~$]4{,}9;\ 5{,}1[$ \quad \textbf{c.}~$n>1000$ \quad \textbf{d.}~$n>6{,}64$
\end{solution}
\begin{fullsolution}
\textbf{a.}~De noemer wordt onbegrensd groot: 0. Hoogste machten: $\tfrac{2x}{x}=2$. De term $-x^2$ beslist: $-\infty$.\\
\textbf{b.}~$-0{,}1<x-5<0{,}1$, dus $4{,}9<x<5{,}1$.\\
\textbf{c.}~Omgekeerde nemen, en dan keert het teken om: $n>1000$.\\
\textbf{d.}~$n>\log_{0{,}5} 0{,}01\approx 6{,}64$. Het teken keert om, want het grondtal is kleiner dan 1.
\end{fullsolution}
\end{oefening}

\begin{oefening}\label{oef:h4-verken}
\blokkop{Verken}\quad Gegeven: $u_n = \dfrac{2n+3}{n}$.
\begin{enumerate}
  \item Bereken $u_{10}$, $u_{100}$ en $u_{1000}$. Welk getal naderen de termen?
  \item Vanaf welk volgnummer ligt $u_n$ minder dan 0,01 van dat getal? En minder dan 0,0001?
  \item Kun je voor elke afstand, hoe klein ook, zo'n volgnummer vinden? Hoe?
\end{enumerate}
\begin{solution}
\textbf{a.}~2,3; 2,03; 2,003; het getal 2 \quad \textbf{b.}~vanaf 301; vanaf 30\,001 \quad \textbf{c.}~ja
\end{solution}
\begin{fullsolution}
\textbf{a.}~$u_n=2+\tfrac3n$, dus 2,3; 2,03 en 2,003. De termen naderen 2.\\
\textbf{b.}~$\tfrac3n<0{,}01$ geeft $n>300$, dus vanaf 301. $\tfrac3n<0{,}0001$ geeft $n>30\,000$, dus vanaf 30\,001.\\
\textbf{c.}~Ja: voor een afstand $\varepsilon>0$ geeft $\tfrac3n<\varepsilon$ dat $n>\tfrac{3}{\varepsilon}$.
\end{fullsolution}
\end{oefening}
\end{oefeningen}

\lkr{Laat oefening \ref{oef:h4-verken} maken vóór de uitleg. Vraag c is de kern van de formele definitie: die komt straks in sectie 2 terug. Stel ook de vraag wat `naderen' eigenlijk betekent.}

%========================================================================================
\section{Convergent en divergent}

\subsection{De begrippen}

\begin{theorie}
Een rij is \emph{convergent} als de termen willekeurig dicht bij een reëel getal $L$ komen en daar blijven. $L$ is de \emph{limiet}: \ $\formule{\displaystyle\lim_{n\to+\infty} u_n = L}$. Een rij die niet convergeert, is \emph{divergent}.

\vspace{1mm}
\noindent\minigrafiek{limiet 2}{-.5}{5.5}{2,4}{1/5,2/3.5,3/3,4/2.75,5/2.6,6/2.5}\hfill
\minigrafiek{limiet $+\infty$}{-.5}{10}{5,10}{1/.5,2/2,3/3,4/5,5/7,6/9.5}\hfill
\minigrafiek{geen limiet}{-3}{3}{-2,2}{1/-2,2/2,3/-2,4/2,5/-2,6/2}
\end{theorie}

De limiet hoeft niet bereikt te worden, maar het mag wel: de rij $5,\ 5,\ 5,\ \ldots$ heeft limiet 5.

\voorbeeld{}{Bepaal de limiet. Is de rij convergent of divergent?\\[.5ex]
\textbf{a.}~$u_n = 4 - \tfrac1n$ \qquad \textbf{b.}~$u_n = 3n - n^2$ \qquad \textbf{c.}~$u_n = 1 + (-1)^n$\label{vb:h4-soort}}{%
\begin{stappen}
\text{a. } \lim_{n\to+\infty} \left(4 - \dfrac1n\right) = 4 & $\tfrac1n$ wordt willekeurig klein: \textbf{convergent}. \\[1.5ex]
\text{b. } \lim_{n\to+\infty} \left(3n - n^2\right) = -\infty & De term $-n^2$ beslist: \textbf{divergent}. \\[1ex]
\text{c. } 0,\ 2,\ 0,\ 2,\ \ldots & De limiet bestaat niet: \textbf{divergent}.
\end{stappen}}

\lkr{Veel leerlingen denken dat een limiet nooit bereikt wordt, en dat divergent hetzelfde is als `naar oneindig gaan'. De constante rij en de rij $1 + (-1)^n$ weerleggen allebei.}

\nujij{Bepaal de limiet van $u_n = \dfrac{(-1)^n}{n}$ en van $u_n = (-1)^n \cdot n$.}{Eerste rij: 0, convergent. Tweede rij: de limiet bestaat niet, divergent.}

\begin{oefeningen}
\begin{oefening}
Bepaal de limiet. Is de rij convergent of divergent?
\begin{reeks}
  \item $u_n = 2 + \tfrac3n$
  \item $u_n = n^2 + 1$
  \item $u_n = 3 \cdot (-1)^n$
\end{reeks}
\begin{solution}
\textbf{a.}~2, convergent \quad \textbf{b.}~$+\infty$, divergent \quad \textbf{c.}~bestaat niet, divergent
\end{solution}
\begin{fullsolution}
\textbf{a.}~$\tfrac3n$ wordt willekeurig klein.\\
\textbf{b.}~$n^2$ groeit onbegrensd.\\
\textbf{c.}~$-3,\ 3,\ {-3},\ \ldots$: de termen naderen geen vast getal.
\end{fullsolution}
\end{oefening}
\end{oefeningen}

\subsection{Rekenkundige en meetkundige rijen}

\begin{theorie}
\begin{tabular}{@{}>{\RaggedRight\arraybackslash}p{6.2cm}>{\RaggedRight\arraybackslash}p{8.4cm}@{}}
\textbf{Rekenkundige rij} & \textbf{Meetkundige rij} ($u_1 \neq 0$) \\[1mm]
$v > 0$: limiet $+\infty$ & $|q| < 1$: limiet 0 \\
$v < 0$: limiet $-\infty$ & $q = 1$: limiet $u_1$ \\
$v = 0$: limiet $u_1$ & $q > 1$: limiet $+\infty$ of $-\infty$, volgens het teken van $u_1$ \\
 & $q \leqslant -1$: de limiet bestaat niet
\end{tabular}
\end{theorie}

\voorbeeld{}{Bepaal de limiet.\\[.5ex]
\textbf{a.}~$9,\ 7,\ 5,\ \ldots$ \qquad \textbf{b.}~$-3,\ {-6},\ {-12},\ \ldots$ \qquad \textbf{c.}~$8,\ {-4},\ 2,\ \ldots$ \qquad \textbf{d.}~$1,\ {-3},\ 9,\ \ldots$\label{vb:h4-rrmr}}{%
\begin{stappen}
\text{a. } v = -2 < 0 & Limiet $-\infty$. \\
\text{b. } q = 2 > 1 \text{ en } u_1 < 0 & Limiet $-\infty$. \\
\text{c. } q = -\tfrac12,\ |q| < 1 & Limiet 0, ook al wisselt het teken. \\
\text{d. } q = -3 \leqslant -1 & De limiet bestaat niet.
\end{stappen}}

\begin{oefeningen}
\begin{oefening}
Bepaal de limiet van \ \textbf{a.}~$2,\ 6,\ 18,\ \ldots$ \quad \textbf{b.}~$100,\ 10,\ 1,\ \ldots$ \quad \textbf{c.}~$5,\ {-5},\ 5,\ \ldots$
\begin{solution}
\textbf{a.}~$+\infty$ \quad \textbf{b.}~0 \quad \textbf{c.}~bestaat niet
\end{solution}
\begin{fullsolution}
\textbf{a.}~$q=3>1$ en $u_1>0$.\\
\textbf{b.}~$q=\tfrac{1}{10}$, dus $|q|<1$.\\
\textbf{c.}~$q=-1$.
\end{fullsolution}
\end{oefening}
\end{oefeningen}

%========================================================================================
\section{De formele definitie}

\begin{theorie}
\begin{minipage}[t]{.57\linewidth}
`Willekeurig dicht' betekent: hoe smal je een band rond $L$ ook kiest, vanaf een bepaald volgnummer liggen alle termen in die band.

\vspace{1.5mm}
$\formule{\displaystyle\lim_{n\to+\infty} u_n = L \ \Leftrightarrow}$

\vspace{.5mm}
$\formule{\forall\, \varepsilon > 0,\ \exists\, N \in \mathbb{N}: \ n > N \Rightarrow |u_n - L| < \varepsilon}$

\vspace{1.5mm}
$\formule{\displaystyle\lim_{n\to+\infty} u_n = +\infty \ \Leftrightarrow}$

\vspace{.5mm}
$\formule{\forall\, r > 0,\ \exists\, N \in \mathbb{N}: \ n > N \Rightarrow u_n > r}$

\vspace{1mm}
Bij $-\infty$ staat er op het einde $u_n < -r$.
\end{minipage}\hfill
\begin{minipage}[t]{.4\linewidth}
\vspace{-2mm}
\begin{tikzpicture}[x=.33cm,y=.72cm,>={Stealth[length=1.6mm]},font=\scriptsize]
  \fill[accent!18] (0,1.4) rectangle (13,2.6);
  \draw[lijngrijs,dashed] (0,2) -- (13,2);
  \draw[->,lijngrijs] (-.1,0) -- (13.5,0) node[right,text=black] {$n$};
  \draw[->,lijngrijs] (0,-.1) -- (0,5.4) node[above,text=black] {$u_n$};
  \draw[lijngrijs,dashed] (5,0) -- (5,2.6);
  \node[below=-1pt] at (5,0) {$N$};
  \node[left=-1pt] at (0,2) {$L$};
  \node[left=-1pt] at (0,2.75) {$L+\varepsilon$};
  \node[left=-1pt] at (0,1.25) {$L-\varepsilon$};
  \foreach \x/\y in {1/5,2/3.5,3/3,4/2.75,5/2.6,6/2.5,7/2.43,8/2.38,9/2.33,10/2.3,11/2.27,12/2.25}
    \fill[accent] (\x,\y) circle (1.5pt);
\end{tikzpicture}
\end{minipage}
\end{theorie}

\thuis{Lees de definitie als een spel. Je tegenstander kiest een $\varepsilon$, hoe klein ook. Jij moet dan een $N$ vinden zodat alle termen na $u_N$ minder dan $\varepsilon$ van $L$ liggen. Lukt dat voor elke $\varepsilon$, dan is $L$ de limiet. De volgorde telt: $N$ mag van $\varepsilon$ afhangen.}

\voorbeeld{}{Bewijs met de definitie: $\displaystyle\lim_{n\to+\infty} \dfrac{2n+3}{n} = 2$.\label{vb:h4-bewijs}}{%
\begin{stappen}
\left|\dfrac{2n+3}{n} - 2\right| = \dfrac{3}{n} & Neem een willekeurige $\varepsilon > 0$ en werk de afstand uit. \\[1.5ex]
\dfrac{3}{n} < \varepsilon \ \Leftrightarrow \ n > \dfrac{3}{\varepsilon} & Los op naar $n$. \\[1.5ex]
N \geqslant \dfrac{3}{\varepsilon},\ N \in \mathbb{N} & Kies zo'n $N$. Dan geldt voor elke $n > N$ dat $\left|u_n - 2\right| < \varepsilon$. \textbf{Bewezen.} \\[1.5ex]
\varepsilon = 0{,}01 \ \Rightarrow \ N = 300 & Controle met oefening \ref{oef:h4-verken}: vanaf $n = 301$.
\end{stappen}}

\lkr{Laat het verschil zien tussen het kladwerk (oplossen naar $n$) en het bewijs zelf (kies $N$, en dan volgt het). Een concreet voorbeeld volstaat volgens het leerplan; de oefeningen op niveau 3 gaan verder.}

\voorbeeld{}{Bewijs met de definitie: $\displaystyle\lim_{n\to+\infty} \sqrt{n} = +\infty$.\label{vb:h4-bewijsinf}}{%
\begin{stappen}
\sqrt{n} > r \ \Leftrightarrow \ n > r^2 & Neem een willekeurige $r > 0$ en los op naar $n$. \\
N \geqslant r^2,\ N \in \mathbb{N} & Kies zo'n $N$. Dan geldt voor elke $n > N$ dat $\sqrt{n} > r$. \textbf{Bewezen.}
\end{stappen}}

\nujij{Bewijs met de definitie: $\displaystyle\lim_{n\to+\infty} \dfrac{5}{n} = 0$.}{$\tfrac5n < \varepsilon$ geeft $n > \tfrac{5}{\varepsilon}$. Kies $N \geqslant \tfrac{5}{\varepsilon}$.}

\begin{oefeningen}
\begin{oefening}
Bewijs met de definitie: $\displaystyle\lim_{n\to+\infty} \dfrac{4n+1}{n} = 4$. Welke $N$ past bij $\varepsilon = 0{,}02$?
\begin{solution}
Kies $N\geqslant\tfrac{1}{\varepsilon}$; bij $\varepsilon=0{,}02$ is $N=50$.
\end{solution}
\begin{fullsolution}
Neem $\varepsilon>0$. $\left|\dfrac{4n+1}{n}-4\right|=\dfrac1n$ en $\dfrac1n<\varepsilon$ geeft $n>\dfrac{1}{\varepsilon}$. Kies $N\geqslant\dfrac{1}{\varepsilon}$: dan geldt voor elke $n>N$ dat $|u_n-4|<\varepsilon$. Bij $\varepsilon=0{,}02$ is $N=50$.
\end{fullsolution}
\end{oefening}
\end{oefeningen}

%========================================================================================
\section{Limieten berekenen}

\subsection{Rekenregels}

\begin{theorie}
Is $u_n = f(n)$ en $\displaystyle\lim_{x\to+\infty} f(x) = L$, dan is ook $\displaystyle\lim_{n\to+\infty} u_n = L$. De rekenregels voor limieten van functies mogen dus: limiet van een som, verschil, product en quotiënt. De rijen $\tfrac1n$ en $\tfrac{1}{n^2}$ hebben limiet 0.
\end{theorie}

Let op: het omgekeerde geldt niet. $u_n = \sin(n\pi)$ is de rij $0,\ 0,\ 0,\ \ldots$ met limiet 0, maar $\displaystyle\lim_{x\to+\infty} \sin(x\pi)$ bestaat niet. En $\infty - \infty$ is niet 0: reken eerst om.

\voorbeeld{}{Bereken \ \textbf{a.}~$\displaystyle\lim_{n\to+\infty} \dfrac{3n^2 - n}{n^2 + 2}$ \quad \textbf{b.}~$\displaystyle\lim_{n\to+\infty} \left(\sqrt{n+1} - \sqrt{n}\right)$\label{vb:h4-reken}}{%
\begin{stappen}
\text{a. } \lim_{n\to+\infty} \dfrac{3 - \tfrac1n}{1 + \tfrac{2}{n^2}} = \dfrac{3 - 0}{1 + 0} = 3 & Deel teller en noemer door de hoogste macht, $n^2$. \\[3ex]
\text{b. } \sqrt{n+1} - \sqrt{n} = \dfrac{1}{\sqrt{n+1} + \sqrt{n}} & Dit is $\infty - \infty$. Vermenigvuldig teller en noemer met de toegevoegde uitdrukking. \\[3ex]
\lim_{n\to+\infty} \dfrac{1}{\sqrt{n+1} + \sqrt{n}} = 0 & De noemer wordt onbegrensd groot.
\end{stappen}}

\nujij{Bereken $\displaystyle\lim_{n\to+\infty} \dfrac{2n - n^2}{4n^2 + 1}$.}{$-\tfrac14$: deel teller en noemer door $n^2$.}

\begin{oefeningen}
\begin{oefening}
Bereken.
\begin{reeks}[2]
  \item $\displaystyle\lim_{n\to+\infty} \dfrac{5n + 2}{3n - 1}$
  \item $\displaystyle\lim_{n\to+\infty} \left(\sqrt{n+4} - \sqrt{n}\right)$
\end{reeks}
\begin{solution}
\textbf{a.}~$\tfrac53$ \quad \textbf{b.}~0
\end{solution}
\begin{fullsolution}
\textbf{a.}~Deel door $n$: $\dfrac{5+\tfrac2n}{3-\tfrac1n}\to\dfrac53$.\\
\textbf{b.}~$\sqrt{n+4}-\sqrt{n}=\dfrac{4}{\sqrt{n+4}+\sqrt{n}}\to 0$.
\end{fullsolution}
\end{oefening}
\end{oefeningen}

\subsection{De insluitstelling}

\begin{theorie}
Als $t_n \leqslant u_n \leqslant v_n$ vanaf een bepaald volgnummer, en $\displaystyle\lim_{n\to+\infty} t_n = \lim_{n\to+\infty} v_n = L$, dan is ook $\displaystyle\lim_{n\to+\infty} u_n = L$.
\end{theorie}

\begin{minipage}[t]{.6\linewidth}
Handig bij een factor die blijft schommelen, zoals $\sin n$, $\cos n$ of $(-1)^n$: die laat zich niet met rekenregels aanpakken, maar is wel begrensd.
\end{minipage}\hfill
\begin{minipage}[t]{.37\linewidth}
\vspace{-1mm}
\begin{tikzpicture}[x=.56cm,y=1.05cm,>={Stealth[length=1.6mm]},font=\scriptsize]
  \draw[->,lijngrijs] (-.1,0) -- (8.8,0) node[right,text=black] {$n$};
  \draw[->,lijngrijs] (0,-1.1) -- (0,1.2) node[above,text=black] {$u_n$};
  \draw[lijngrijs,dashed,domain=.9:8.5,samples=40] plot (\x,{1/\x});
  \draw[lijngrijs,dashed,domain=.9:8.5,samples=40] plot (\x,{-1/\x});
  \node[right=-1pt] at (1.25,.82) {$\tfrac1n$};
  \node[right=-1pt] at (1.25,-.82) {$-\tfrac1n$};
  \foreach \x/\y in {1/.54,2/-.208,3/-.33,4/-.163,5/.057,6/.16,7/.108,8/-.018}
    \fill[accent] (\x,\y) circle (1.6pt);
\end{tikzpicture}
\end{minipage}

\voorbeeld{}{Bereken $\displaystyle\lim_{n\to+\infty} \dfrac{\cos n}{n}$. (Zie de figuur.)\label{vb:h4-insluit}}{%
\begin{stappen}
-1 \leqslant \cos n \leqslant 1 & De cosinus is begrensd. \\[1ex]
-\dfrac1n \leqslant \dfrac{\cos n}{n} \leqslant \dfrac1n & Delen door $n > 0$. \\[1.5ex]
\lim_{n\to+\infty} \dfrac{\cos n}{n} = 0 & Beide grenzen hebben limiet 0: de rij wordt ingesloten.
\end{stappen}}

\begin{oefeningen}
\begin{oefening}
Bereken $\displaystyle\lim_{n\to+\infty} \dfrac{\sin n}{n^2}$.
\begin{solution}
0
\end{solution}
\begin{fullsolution}
$-\dfrac{1}{n^2}\leqslant\dfrac{\sin n}{n^2}\leqslant\dfrac{1}{n^2}$, en beide grenzen hebben limiet 0.
\end{fullsolution}
\end{oefening}
\end{oefeningen}

%========================================================================================
\section{Stijgende en dalende rijen}

\begin{theorie}
Een rij is \emph{stijgend} als $u_{n+1} > u_n$ voor elke $n$, en \emph{dalend} als $u_{n+1} < u_n$ voor elke $n$. Een rij is \emph{naar boven begrensd} als er een $M$ is met $u_n \leqslant M$ voor elke $n$, en \emph{naar onder begrensd} als er een $m$ is met $u_n \geqslant m$ voor elke $n$. Begrensd is allebei.

\vspace{1.5mm}
\textbf{Eigenschap.} Een stijgende, naar boven begrensde rij is convergent. Een dalende, naar onder begrensde rij ook.
\end{theorie}

Eén voorwaarde volstaat niet: $u_n = n$ is stijgend en $u_n = (-1)^n$ is begrensd, maar beide zijn divergent.

\thuis{Stel je de grafiek voor: de punten gaan alleen maar omhoog, maar er is een plafond dat ze niet kunnen doorbreken. Dan moeten ze zich wel ophopen onder een bepaalde hoogte. Die hoogte is de limiet. Het bewijs gebruiken we hier niet.}

\voorbeeld{}{Toon aan dat $u_n = \dfrac{n}{n+2}$ stijgend en naar boven begrensd is. Wat is de limiet?\label{vb:h4-monotoon}}{%
\begin{stappen}
u_n = 1 - \dfrac{2}{n+2} & Herschrijf. Dan is $u_n < 1$ voor elke $n$: bovengrens 1. \\[2ex]
u_{n+1} - u_n = \dfrac{2}{n+2} - \dfrac{2}{n+3} & Bereken het verschil van twee opeenvolgende termen. \\[2ex]
\phantom{u_{n+1} - u_n} = \dfrac{2}{(n+2)(n+3)} > 0 & Positief voor elke $n$: stijgend, dus convergent. \\[2ex]
\lim_{n\to+\infty} u_n = 1 & De eigenschap zegt dat er een limiet is. De rekenregels geven hier de waarde.
\end{stappen}}

\begin{oefeningen}
\begin{oefening}
Toon aan dat $u_n = 4 + \dfrac{3}{n}$ dalend en naar onder begrensd is. Wat is de limiet?
\begin{solution}
Ondergrens 4; limiet 4
\end{solution}
\begin{fullsolution}
$u_n>4$ voor elke $n$, dus 4 is een ondergrens. $u_{n+1}-u_n=\dfrac{3}{n+1}-\dfrac3n=\dfrac{-3}{n(n+1)}<0$, dus de rij is dalend. De limiet is 4.
\end{fullsolution}
\end{oefening}
\end{oefeningen}

%========================================================================================
\section{Discrete dynamische processen}

\begin{theorie}
Een proces met $u_n = a\,u_{n-1} + b$ en $|a| < 1$ gaat naar een \emph{evenwicht}. Dat vind je door $u_n$ en $u_{n-1}$ gelijk te stellen aan de limiet $L$:

\vspace{1mm}
\qquad $\formule{L = a L + b}$
\end{theorie}

\begin{minipage}[t]{.6\linewidth}
Discreet: je kijkt met vaste tussenpozen. Dynamisch: er verandert iets. In het evenwicht blijft de hoeveelheid gelijk, maar er verdwijnt en er komt evenveel bij.
\end{minipage}\hfill
\begin{minipage}[t]{.37\linewidth}
\vspace{-1mm}
\begin{tikzpicture}[x=.39cm,y=.0095cm,>={Stealth[length=1.6mm]},font=\scriptsize]
  \draw[->,lijngrijs] (-.1,0) -- (11,0) node[right,text=black] {$n$};
  \draw[->,lijngrijs] (0,-10) -- (0,450) node[above,text=black] {$u_n$};
  \foreach \x in {2,4,...,10} \draw[lijngrijs] (\x,6) -- (\x,-6) node[below=-1pt,text=grijs] {\x};
  \foreach \y in {100,200,300,400} \draw[lijngrijs] (.07,\y) -- (-.07,\y) node[left=-1pt,text=grijs] {\y};
  \draw[lijngrijs,dashed] (0,250) -- (10.6,250) node[below=-1pt,pos=.9,text=black] {250};
  \foreach \x/\y in {0/400,1/370,2/346,3/326.8,4/311.4,5/299.2,6/289.3,7/281.5,8/275.2,9/270.1,10/266.1}
    \fill[accent] (\x,\y) circle (1.6pt);
\end{tikzpicture}
\end{minipage}

\voorbeeld{}{Het zwembad uit hoofdstuk 1: $u_0 = 400$ en $u_n = 0{,}8\,u_{n-1} + 50$. Bepaal het evenwicht $L$. Toon aan dat de rij $v_n = u_n - L$ meetkundig is, en leid daaruit af dat de limiet echt $L$ is.\label{vb:h4-dyn}}{%
\begin{stappen}
L = 0{,}8\,L + 50 \ \Rightarrow \ L = 250 & Het evenwicht: 250 g chloor. \\
v_n = u_n - 250 = 0{,}8\,u_{n-1} - 200 & Vervang $u_n$ door het voorschrift. \\
\phantom{v_n} = 0{,}8\,(u_{n-1} - 250) = 0{,}8\,v_{n-1} & Meetkundig met $q = 0{,}8$ en $v_0 = 150$. \\
u_n = 250 + 150 \cdot 0{,}8^{\,n} & Controle: $u_1 = 250 + 120 = 370$. \\
\lim_{n\to+\infty} u_n = 250 & Want $|q| < 1$, dus $0{,}8^{\,n} \to 0$.
\end{stappen}}

\lkr{De vergelijking $L = aL + b$ veronderstelt dat de limiet bestaat. De hulprij $v_n$ bewijst dat, en sluit aan bij de hulprij uit hoofdstuk 1. Bij $|a| \geqslant 1$ levert de vergelijking vaak nog een getal op, ook al is er geen limiet: zie de opdrachten.}

\nujij{$u_0 = 90$ en $u_n = 0{,}5\,u_{n-1} + 10$. Bepaal het evenwicht en een expliciet voorschrift.}{$L = 20$ en $u_n = 20 + 70 \cdot 0{,}5^{\,n}$.}

\begin{oefeningen}
\begin{oefening}
Een bedrijf heeft 300 werknemers. Elk jaar vertrekt 10\% en komen er 20 nieuwe bij. Geef een recursief voorschrift. Hoeveel werknemers zijn er op lange termijn?
\begin{solution}
$u_n=0{,}9\,u_{n-1}+20$, $u_0=300$; 200
\end{solution}
\begin{fullsolution}
Na het vertrek blijft 90\% over, dus $u_n=0{,}9\,u_{n-1}+20$ met $u_0=300$. $L=0{,}9\,L+20$ geeft $L=200$.
\end{fullsolution}
\end{oefening}
\end{oefeningen}

%========================================================================================
\section{De meetkundige reeks}

\begin{theorie}
Een \emph{reeks} is een oneindige som. Haar waarde is de limiet van de sommen $s_n$. Bij een meetkundige rij met $|q| < 1$ gaat $q^{\,n}$ naar 0, dus:

\vspace{1mm}
\qquad $u_1 + u_1 q + u_1 q^2 + \cdots = \displaystyle\lim_{n\to+\infty} u_1 \cdot \frac{q^{\,n} - 1}{q - 1} = \formule{\dfrac{u_1}{1 - q}}$

\vspace{1.5mm}
Bij $|q| \geqslant 1$ en $u_1 \neq 0$ heeft de reeks geen eindige waarde: ze \emph{divergeert}.
\end{theorie}

\voorbeeld{}{Bereken $8 + 2 + \tfrac12 + \cdots$ en schrijf $0{,}272727\ldots$ als een breuk.\label{vb:h4-reeks}}{%
\begin{stappen}
q = \dfrac14,\quad \dfrac{8}{1 - \tfrac14} = \dfrac{32}{3} & $|q| < 1$: de formule mag. \\[2ex]
0{,}27 + 0{,}0027 + 0{,}000027 + \cdots & Schrijf het getal als een som: $u_1 = 0{,}27$ en $q = 0{,}01$. \\[1ex]
\dfrac{0{,}27}{1 - 0{,}01} = \dfrac{27}{99} = \dfrac{3}{11} & Controle met het rekentoestel: $3 : 11 = 0{,}2727\ldots$
\end{stappen}}

\nujij{Bereken $9 - 3 + 1 - \cdots$ en schrijf $0{,}1515\ldots$ als een breuk.}{$\tfrac{9}{1 + 1/3} = \tfrac{27}{4}$ en $\tfrac{15}{99} = \tfrac{5}{33}$.}

\lkr{De slak uit hoofdstuk 3 past hier: $1 + \tfrac12 + \tfrac14 + \cdots = 2$. De reeks is een suggestie in het leerplan, geen doel.}

\begin{oefeningen}
\begin{oefening}
Bereken $12 + 4 + \tfrac43 + \cdots$ en schrijf $0{,}888\ldots$ als een breuk.
\begin{solution}
18 en $\tfrac89$
\end{solution}
\begin{fullsolution}
$q=\tfrac13$, dus $\dfrac{12}{1-\tfrac13}=18$. Verder is $0{,}8+0{,}08+\cdots=\dfrac{0{,}8}{1-0{,}1}=\dfrac89$.
\end{fullsolution}
\end{oefening}
\end{oefeningen}

%----------------------------------------------------------------------------------------
% OPDRACHTEN: door elkaar. Niveau 1, dan 2, dan 3.
%----------------------------------------------------------------------------------------
\opdrachtenkop

\lkr{De opdrachten staan bewust door elkaar, met stijgende moeilijkheid. In deze versie staat bij elke opdracht het onderwerp, zodat je kunt kiezen. Vraag bij elke limiet ook: convergent of divergent?}

\begin{oefeningen}
\begin{oefening}[1]\label{oef:h4-soort}
\thema{convergent en divergent}Bepaal de limiet. Is de rij convergent of divergent?
\begin{reeks}
  \item $u_n = 5 + \tfrac2n$
  \item $u_n = n^2 - 3$
  \item $u_n = 4 \cdot (-1)^n$
  \item $u_n = 7$
  \item $u_n = \tfrac{(-1)^n}{n}$
  \item $u_n = 1 - n$
\end{reeks}
\begin{solution}
\textbf{a.}~5, conv. \quad \textbf{b.}~$+\infty$, div. \quad \textbf{c.}~bestaat niet, div. \quad \textbf{d.}~7, conv. \quad \textbf{e.}~0, conv. \quad \textbf{f.}~$-\infty$, div.
\end{solution}
\begin{fullsolution}
\textbf{a.}~$\tfrac2n\to 0$.\\
\textbf{b.}~$n^2$ groeit onbegrensd.\\
\textbf{c.}~$-4,\ 4,\ {-4},\ \ldots$: geen vast getal.\\
\textbf{d.}~Een constante rij heeft die constante als limiet.\\
\textbf{e.}~Het teken wisselt, maar de termen komen willekeurig dicht bij 0.\\
\textbf{f.}~De termen dalen onbegrensd.
\end{fullsolution}
\end{oefening}

\begin{oefening}[1]\label{oef:h4-reeks}
\thema{meetkundige reeks}Bereken, of zeg dat de reeks divergeert.
\begin{reeks}[2]
  \item $6 + 3 + \tfrac32 + \cdots$
  \item $1 - \tfrac12 + \tfrac14 - \cdots$
  \item $1 + 3 + 9 + \cdots$
  \item $0{,}444\ldots$
\end{reeks}
\begin{solution}
\textbf{a.}~12 \quad \textbf{b.}~$\tfrac23$ \quad \textbf{c.}~divergeert \quad \textbf{d.}~$\tfrac49$
\end{solution}
\begin{fullsolution}
\textbf{a.}~$q=\tfrac12$: $\dfrac{6}{1-\tfrac12}=12$.\\
\textbf{b.}~$q=-\tfrac12$: $\dfrac{1}{1+\tfrac12}=\dfrac23$.\\
\textbf{c.}~$q=3$: de formule mag niet.\\
\textbf{d.}~$u_1=0{,}4$ en $q=0{,}1$: $\dfrac{0{,}4}{0{,}9}=\dfrac49$.
\end{fullsolution}
\end{oefening}

\begin{oefening}[1]\label{oef:h4-ratio}
\thema{rekenregels}Bereken.
\begin{reeks}[2]
  \item $\displaystyle\lim_{n\to+\infty} \dfrac{5n - 1}{2n + 3}$
  \item $\displaystyle\lim_{n\to+\infty} \dfrac{n^2 + 1}{n^3 - 2}$
  \item $\displaystyle\lim_{n\to+\infty} \dfrac{4 - n^2}{n^2 + n}$
  \item $\displaystyle\lim_{n\to+\infty} \dfrac{n^2 + 3}{n + 1}$
\end{reeks}
\begin{solution}
\textbf{a.}~$\tfrac52$ \quad \textbf{b.}~0 \quad \textbf{c.}~$-1$ \quad \textbf{d.}~$+\infty$
\end{solution}
\begin{fullsolution}
Vergelijk de hoogste machten in teller en noemer.\\
\textbf{a.}~$\dfrac{5n}{2n}=\dfrac52$.\\
\textbf{b.}~De noemer heeft de hoogste macht: 0.\\
\textbf{c.}~$\dfrac{-n^2}{n^2}=-1$.\\
\textbf{d.}~De teller heeft de hoogste macht, en beide zijn positief: $+\infty$.
\end{fullsolution}
\end{oefening}

\begin{oefening}[1]\label{oef:h4-rrmr}
\thema{rekenkundige en meetkundige rijen}Bepaal de limiet.
\begin{reeks}
  \item $2,\ 5,\ 8,\ \ldots$
  \item $100,\ 50,\ 25,\ \ldots$
  \item $3,\ {-6},\ 12,\ \ldots$
  \item $-1,\ {-3},\ {-9},\ \ldots$
  \item $9,\ 7,\ 5,\ \ldots$
  \item $8,\ {-4},\ 2,\ \ldots$
\end{reeks}
\begin{solution}
\textbf{a.}~$+\infty$ \quad \textbf{b.}~0 \quad \textbf{c.}~bestaat niet \quad \textbf{d.}~$-\infty$ \quad \textbf{e.}~$-\infty$ \quad \textbf{f.}~0
\end{solution}
\begin{fullsolution}
\textbf{a.}~$v=3>0$.\\
\textbf{b.}~$q=\tfrac12$.\\
\textbf{c.}~$q=-2$.\\
\textbf{d.}~$q=3$ en $u_1<0$.\\
\textbf{e.}~$v=-2<0$.\\
\textbf{f.}~$q=-\tfrac12$.
\end{fullsolution}
\end{oefening}

\begin{oefening}[1]
\thema{fouten: limiet bereiken}Iemand schrijft: `Een limiet wordt nooit bereikt, dus de rij $5,\ 5,\ 5,\ \ldots$ heeft geen limiet.' Wat is er fout?
\begin{solution}
De limiet is 5: de limiet mag bereikt worden.
\end{solution}
\begin{fullsolution}
Voor elke $\varepsilon>0$ is $|u_n-5|=0<\varepsilon$, dus de limiet is 5. De definitie vraagt niet dat de termen van $L$ verschillen.
\end{fullsolution}
\end{oefening}

\begin{oefening}[1]\label{oef:h4-monotoon}
\thema{stijgend, dalend, begrensd}Is de rij stijgend, dalend of geen van beide? Is ze naar boven begrensd? Naar onder? Is ze convergent?
\begin{reeks}[4]
  \item $u_n = \tfrac1n$
  \item $u_n = n^2$
  \item $u_n = (-1)^n$
  \item $u_n = 5 - \tfrac1n$
\end{reeks}
\begin{solution}
\textbf{a.}~dalend; boven 1, onder 0; conv. \quad \textbf{b.}~stijgend; enkel onder (1); div. \quad \textbf{c.}~geen van beide; boven 1, onder $-1$; div. \quad \textbf{d.}~stijgend; boven 5, onder 4; conv.
\end{solution}
\begin{fullsolution}
\textbf{a.}~Dalend, met $0<u_n\leqslant 1$. Limiet 0.\\
\textbf{b.}~Stijgend, met $u_n\geqslant 1$, en geen bovengrens. Limiet $+\infty$.\\
\textbf{c.}~Wisselt, met $-1\leqslant u_n\leqslant 1$. Begrensd, maar de limiet bestaat niet.\\
\textbf{d.}~Stijgend, met $4\leqslant u_n<5$. Limiet 5.
\end{fullsolution}
\end{oefening}

\begin{oefening}[1]\label{oef:h4-evenwicht}
\thema{dynamisch proces}Bepaal het evenwicht.
\begin{reeks}
  \item $u_n = 0{,}5\,u_{n-1} + 6$
  \item $u_n = 0{,}9\,u_{n-1} + 20$
  \item $u_n = 0{,}75\,u_{n-1} - 3$
\end{reeks}
\begin{solution}
\textbf{a.}~12 \quad \textbf{b.}~200 \quad \textbf{c.}~$-12$
\end{solution}
\begin{fullsolution}
\textbf{a.}~$0{,}5\,L=6$.\\
\textbf{b.}~$0{,}1\,L=20$.\\
\textbf{c.}~$0{,}25\,L=-3$.
\end{fullsolution}
\end{oefening}

\begin{oefening}[1]\label{oef:h4-epsgetal}
\thema{definitie: met getallen}Vanaf welk volgnummer geldt dit?
\begin{enumerate}
  \item $u_n = \dfrac{1}{n^2}$ ligt minder dan 0,0001 van 0.
  \item $u_n = \dfrac{n+1}{n}$ ligt minder dan 0,001 van 1.
  \item $u_n = \dfrac{(-1)^n}{n}$ ligt minder dan 0,05 van 0.
\end{enumerate}
\begin{solution}
\textbf{a.}~101 \quad \textbf{b.}~1001 \quad \textbf{c.}~21
\end{solution}
\begin{fullsolution}
\textbf{a.}~$\dfrac{1}{n^2}<0{,}0001$ geeft $n^2>10\,000$, dus $n>100$.\\
\textbf{b.}~$\left|\dfrac{n+1}{n}-1\right|=\dfrac1n<0{,}001$ geeft $n>1000$.\\
\textbf{c.}~$\left|\dfrac{(-1)^n}{n}\right|=\dfrac1n<0{,}05$ geeft $n>20$.
\end{fullsolution}
\end{oefening}

\begin{oefening}[2]\label{oef:h4-bewijs}
\thema{definitie: bewijs}Bewijs met de definitie.
\begin{reeks}
  \item $\displaystyle\lim_{n\to+\infty} \dfrac{7}{n} = 0$
  \item $\displaystyle\lim_{n\to+\infty} \dfrac{4n+1}{2n} = 2$
  \item $\displaystyle\lim_{n\to+\infty} \dfrac{n-1}{n+3} = 1$
\end{reeks}
\begin{solution}
\textbf{a.}~$N\geqslant\tfrac{7}{\varepsilon}$ \quad \textbf{b.}~$N\geqslant\tfrac{1}{2\varepsilon}$ \quad \textbf{c.}~$N\geqslant\tfrac{4}{\varepsilon}$
\end{solution}
\begin{fullsolution}
Neem telkens een willekeurige $\varepsilon>0$.\\
\textbf{a.}~$\tfrac7n<\varepsilon$ geeft $n>\tfrac{7}{\varepsilon}$. Kies $N\geqslant\tfrac{7}{\varepsilon}$.\\
\textbf{b.}~$\left|\dfrac{4n+1}{2n}-2\right|=\dfrac{1}{2n}<\varepsilon$ geeft $n>\dfrac{1}{2\varepsilon}$. Kies $N\geqslant\dfrac{1}{2\varepsilon}$.\\
\textbf{c.}~$\left|\dfrac{n-1}{n+3}-1\right|=\dfrac{4}{n+3}<\varepsilon$ geeft $n>\dfrac{4}{\varepsilon}-3$. Elke $N\geqslant\dfrac{4}{\varepsilon}$ is goed, want dan is ook $n>\dfrac{4}{\varepsilon}-3$.
\end{fullsolution}
\end{oefening}

\begin{oefening}[2]
\thema{context: dynamisch proces}Een sportclub telt 800 leden. Elk jaar stopt 25\% en sluiten er 150 nieuwe leden aan.
\begin{enumerate}
  \item Geef een recursief voorschrift. Hoeveel leden zijn er op lange termijn?
  \item Toon aan dat $u_n = 600 + 200 \cdot 0{,}75^{\,n}$.
  \item Na hoeveel jaar zijn er voor het eerst minder dan 650 leden?
  \item Hoeveel nieuwe leden per jaar zijn nodig om op lange termijn 1000 leden te hebben?
\end{enumerate}
\begin{solution}
\textbf{a.}~$u_n=0{,}75\,u_{n-1}+150$, $u_0=800$; 600 \quad \textbf{c.}~na 5 jaar \quad \textbf{d.}~250
\end{solution}
\begin{fullsolution}
\textbf{a.}~$L=0{,}75\,L+150$ geeft $L=600$.\\
\textbf{b.}~$v_n=u_n-600=0{,}75\,u_{n-1}-450=0{,}75\,v_{n-1}$ met $v_0=200$, dus $v_n=200\cdot 0{,}75^{\,n}$.\\
\textbf{c.}~$200\cdot 0{,}75^{\,n}<50$ geeft $0{,}75^{\,n}<0{,}25$ en $n>4{,}82$, dus na 5 jaar.\\
\textbf{d.}~$1000=750+b$ geeft $b=250$.
\end{fullsolution}
\end{oefening}

\begin{oefening}[2]
\thema{definitie: oneindig}Bewijs met de definitie.
\begin{reeks}[2]
  \item $\displaystyle\lim_{n\to+\infty} (2n + 1) = +\infty$
  \item $\displaystyle\lim_{n\to+\infty} (5 - n^2) = -\infty$
\end{reeks}
\begin{solution}
\textbf{a.}~$N\geqslant\tfrac{r-1}{2}$ \quad \textbf{b.}~$N\geqslant\sqrt{r+5}$
\end{solution}
\begin{fullsolution}
\textbf{a.}~Neem $r>0$. $2n+1>r$ geeft $n>\dfrac{r-1}{2}$. Kies $N\in\mathbb{N}$ met $N\geqslant\dfrac{r-1}{2}$.\\
\textbf{b.}~Neem $r>0$. $5-n^2<-r$ geeft $n^2>r+5$, dus $n>\sqrt{r+5}$. Kies $N\in\mathbb{N}$ met $N\geqslant\sqrt{r+5}$.
\end{fullsolution}
\end{oefening}

\begin{oefening}[2]\label{oef:h4-insluit}
\thema{insluitstelling}Bereken.
\begin{reeks}
  \item $\displaystyle\lim_{n\to+\infty} \dfrac{(-1)^n}{\sqrt{n}}$
  \item $\displaystyle\lim_{n\to+\infty} \dfrac{3 + \cos n}{n}$
  \item $\displaystyle\lim_{n\to+\infty} \left(2 + \dfrac{\sin n}{n}\right)$
\end{reeks}
\begin{solution}
\textbf{a.}~0 \quad \textbf{b.}~0 \quad \textbf{c.}~2
\end{solution}
\begin{fullsolution}
\textbf{a.}~$-\dfrac{1}{\sqrt n}\leqslant u_n\leqslant\dfrac{1}{\sqrt n}$, en beide grenzen gaan naar 0.\\
\textbf{b.}~$2\leqslant 3+\cos n\leqslant 4$, dus $\dfrac2n\leqslant u_n\leqslant\dfrac4n$.\\
\textbf{c.}~$2-\dfrac1n\leqslant u_n\leqslant 2+\dfrac1n$, en beide grenzen gaan naar 2.
\end{fullsolution}
\end{oefening}

\begin{oefening}[2]
\thema{context: reeks}Een bal valt van 3 m hoogte en stuit telkens terug tot 80\% van de vorige hoogte. Welke afstand legt de bal in totaal af, als hij blijft stuiten tot hij stil ligt?
\begin{solution}
27 m
\end{solution}
\begin{fullsolution}
De stuithoogtes $2{,}4;\ 1{,}92;\ \ldots$ tel je dubbel, omhoog en omlaag. De som is $\dfrac{2{,}4}{1-0{,}8}=12$, dus de afstand is $3+2\cdot 12=27$ m.
\end{fullsolution}
\end{oefening}

\begin{oefening}[2]
\thema{rekenregels: machten}Bereken.
\begin{reeks}
  \item $\displaystyle\lim_{n\to+\infty} \dfrac{3^n + 1}{3^n}$
  \item $\displaystyle\lim_{n\to+\infty} \dfrac{2^n - 1}{3^n}$
  \item $\displaystyle\lim_{n\to+\infty} \dfrac{2^{n+1} + 3^n}{3^{n+1}}$
\end{reeks}
\begin{solution}
\textbf{a.}~1 \quad \textbf{b.}~0 \quad \textbf{c.}~$\tfrac13$
\end{solution}
\begin{fullsolution}
\textbf{a.}~$1+(\tfrac13)^n\to 1$.\\
\textbf{b.}~$(\tfrac23)^n-(\tfrac13)^n\to 0$.\\
\textbf{c.}~$\tfrac23\cdot(\tfrac23)^n+\tfrac13\to\tfrac13$.
\end{fullsolution}
\end{oefening}

\begin{oefening}[2]
\thema{functie en rij}Gegeven: $u_n = \sin(2n\pi)$ en $f(x) = \sin(2x\pi)$.
\begin{enumerate}
  \item Bereken $u_1$, $u_2$ en $u_3$. Wat is de limiet van de rij?
  \item Bestaat $\displaystyle\lim_{x\to+\infty} f(x)$? Wat leer je hieruit?
\end{enumerate}
\begin{solution}
\textbf{a.}~0, 0, 0; limiet 0 \quad \textbf{b.}~Nee.
\end{solution}
\begin{fullsolution}
\textbf{a.}~$\sin(2n\pi)=0$ voor elk natuurlijk getal $n$, dus de limiet is 0.\\
\textbf{b.}~$f$ blijft schommelen tussen $-1$ en 1, dus die limiet bestaat niet. Uit de limiet van de rij volgt dus niets over de limiet van de functie: alleen het omgekeerde geldt.
\end{fullsolution}
\end{oefening}

\begin{oefening}[2]
\thema{dynamisch proces: niet stijgend of dalend}Gegeven: $u_0 = 10$ en $u_n = -0{,}5\,u_{n-1} + 6$.
\begin{enumerate}
  \item Bereken $u_1$, $u_2$ en $u_3$. Is de rij stijgend of dalend?
  \item Bepaal het evenwicht $L$ en een expliciet voorschrift.
  \item De rij is niet stijgend of dalend en toch convergent. Is dat in strijd met de eigenschap uit sectie 4?
\end{enumerate}
\begin{solution}
\textbf{a.}~1; 5,5; 3,25; geen van beide \quad \textbf{b.}~$L=4$, $u_n=4+6\cdot(-0{,}5)^n$ \quad \textbf{c.}~nee
\end{solution}
\begin{fullsolution}
\textbf{a.}~$-5+6=1$, $-0{,}5+6=5{,}5$, $-2{,}75+6=3{,}25$. De termen gaan op en neer.\\
\textbf{b.}~$L=-0{,}5\,L+6$ geeft $L=4$. Met $v_n=u_n-4$ is $v_n=-0{,}5\,v_{n-1}$ en $v_0=6$.\\
\textbf{c.}~Nee: de eigenschap zegt dat stijgend en begrensd voldoende is om te convergeren, niet dat het nodig is.
\end{fullsolution}
\end{oefening}

\begin{oefening}[2]
\thema{reeks: repeterende decimalen}Schrijf als een breuk.
\begin{reeks}
  \item $0{,}121212\ldots$
  \item $1{,}2555\ldots$
  \item $0{,}999\ldots$
\end{reeks}
\begin{solution}
\textbf{a.}~$\tfrac{4}{33}$ \quad \textbf{b.}~$\tfrac{113}{90}$ \quad \textbf{c.}~1
\end{solution}
\begin{fullsolution}
\textbf{a.}~$\dfrac{0{,}12}{1-0{,}01}=\dfrac{12}{99}=\dfrac{4}{33}$.\\
\textbf{b.}~$1{,}2+\dfrac{0{,}05}{1-0{,}1}=\dfrac65+\dfrac{1}{18}=\dfrac{113}{90}$.\\
\textbf{c.}~$\dfrac{0{,}9}{1-0{,}1}=1$. Dus $0{,}999\ldots$ en 1 zijn hetzelfde getal.
\end{fullsolution}
\end{oefening}

\begin{oefening}[2]
\thema{rekenregels: toegevoegde}Bereken.
\begin{reeks}[2]
  \item $\displaystyle\lim_{n\to+\infty} \left(\sqrt{n^2 + n} - n\right)$
  \item $\displaystyle\lim_{n\to+\infty} \left(n^2 - n\right)$
\end{reeks}
\begin{solution}
\textbf{a.}~$\tfrac12$ \quad \textbf{b.}~$+\infty$
\end{solution}
\begin{fullsolution}
\textbf{a.}~$\sqrt{n^2+n}-n=\dfrac{n}{\sqrt{n^2+n}+n}=\dfrac{1}{\sqrt{1+\tfrac1n}+1}\to\dfrac12$.\\
\textbf{b.}~$n(n-1)$: een product van twee factoren die onbegrensd groeien. Het antwoord 0, uit $\infty-\infty$, is fout.
\end{fullsolution}
\end{oefening}

\begin{oefening}[2]
\begin{metprent}{schildpad}
\thema{context: reeks}Een loper start 100 m achter een schildpad. De loper loopt 10 m/s, de schildpad 1 m/s. Als de loper op de startplaats van de schildpad is, is die 10 m verder. Als de loper daar is, is ze nog 1 m verder, enzovoort.
\begin{enumerate}
  \item Schrijf de tijd tot de loper de schildpad inhaalt als een oneindige som. Bereken ze.
  \item Controleer met een vergelijking.
\end{enumerate}
\end{metprent}
\begin{solution}
\textbf{a.}~$10+1+0{,}1+\cdots=\tfrac{100}{9}$ s, ongeveer 11,11 s
\end{solution}
\begin{fullsolution}
\textbf{a.}~De stukken duren 10 s, 1 s, 0,1 s, \ldots, dus $\dfrac{10}{1-0{,}1}=\dfrac{100}{9}$ s.\\
\textbf{b.}~$10t=100+t$ geeft $t=\dfrac{100}{9}$. Een oneindig aantal stukken kan dus samen een eindige tijd duren.
\end{fullsolution}
\end{oefening}

\begin{oefening}[2]
\thema{reeks: parameter}Voor welke $x \neq 1$ convergeert $1 + (x-1) + (x-1)^2 + \cdots$? Wat is dan de som?
\begin{solution}
$0<x<2$ met $x\neq 1$; de som is $\dfrac{1}{2-x}$
\end{solution}
\begin{fullsolution}
$q=x-1$, en $|x-1|<1$ geeft $0<x<2$. De som is $\dfrac{1}{1-(x-1)}=\dfrac{1}{2-x}$.
\end{fullsolution}
\end{oefening}

\begin{oefening}[2]
\thema{rekenregels: parameter}Bepaal $k$ zodat $\displaystyle\lim_{n\to+\infty} \dfrac{k\,n^2 - 1}{4n^2 + n} = 3$.
\begin{solution}
$k=12$
\end{solution}
\begin{fullsolution}
De hoogste machten geven $\dfrac{k}{4}=3$, dus $k=12$.
\end{fullsolution}
\end{oefening}

\begin{oefening}[3]
\thema{dynamisch proces: voorwaarde}Bekijk $u_n = a\,u_{n-1} + b$ met $a \neq 1$.
\begin{enumerate}
  \item Stel $L = \dfrac{b}{1-a}$ en $v_n = u_n - L$. Toon aan dat $v_n = a\,v_{n-1}$.
  \item Voor welke $a$ is de rij convergent, als $u_0 \neq L$?
  \item Neem $u_0 = 1$ en $u_n = 2u_{n-1} - 3$. De vergelijking $L = 2L - 3$ geeft $L = 3$. Is 3 de limiet?
\end{enumerate}
\begin{solution}
\textbf{b.}~$|a|<1$ \quad \textbf{c.}~nee: limiet $-\infty$
\end{solution}
\begin{fullsolution}
\textbf{a.}~Uit $L=aL+b$ volgt $v_n=au_{n-1}+b-L=au_{n-1}-aL=a\,v_{n-1}$.\\
\textbf{b.}~$v_n=v_0\cdot a^n$ met $v_0\neq 0$. Dat gaat enkel naar een eindige limiet als $|a|<1$ (want $a\neq 1$).\\
\textbf{c.}~Nee: $1,\ {-1},\ {-5},\ {-13},\ \ldots$ Hier is $u_n=3-2\cdot 2^n$, met limiet $-\infty$. De vergelijking vindt alleen de waarde die de limiet zou moeten hebben, als hij bestaat.
\end{fullsolution}
\end{oefening}

\begin{oefening}[3]
\thema{stijgend en begrensd: inductie}Gegeven: $u_1 = 1$ en $u_n = \sqrt{2 + u_{n-1}}$.
\begin{enumerate}
  \item Bereken $u_2$, $u_3$ en $u_4$ op drie decimalen.
  \item Bewijs met volledige inductie dat $0 < u_n < 2$ voor elke $n$.
  \item Toon aan dat de rij stijgend is. Kijk daarvoor naar $u_{n+1}^{\,2} - u_n^{\,2}$.
  \item Waarom is de rij convergent? Bereken de limiet.
\end{enumerate}
\begin{solution}
\textbf{a.}~1,732; 1,932; 1,983 \quad \textbf{d.}~2
\end{solution}
\begin{fullsolution}
\textbf{a.}~$\sqrt3\approx 1{,}732$, $\sqrt{3{,}732}\approx 1{,}932$, $\sqrt{3{,}932}\approx 1{,}983$.\\
\textbf{b.}~Basis: $0<1<2$. Stap: uit $0<u_k<2$ volgt $0<\sqrt{2+u_k}<\sqrt4=2$.\\
\textbf{c.}~$u_{n+1}^{\,2}-u_n^{\,2}=2+u_n-u_n^{\,2}=(2-u_n)(1+u_n)>0$, en alle termen zijn positief.\\
\textbf{d.}~Stijgend en naar boven begrensd. Voor de limiet geldt $L=\sqrt{2+L}$, dus $L^2-L-2=0$ en $L=2$ of $L=-1$. Omdat de termen positief zijn, is $L=2$.
\end{fullsolution}
\end{oefening}

\begin{oefening}[3]
\thema{definitie: rekenregel}Stel dat $\displaystyle\lim_{n\to+\infty} u_n = L$ en $k \neq 0$. Bewijs met de definitie dat $\displaystyle\lim_{n\to+\infty} k\,u_n = kL$.
\begin{solution}
Gebruik de $N$ die bij $\tfrac{\varepsilon}{|k|}$ hoort.
\end{solution}
\begin{fullsolution}
Neem $\varepsilon>0$. Dan is ook $\dfrac{\varepsilon}{|k|}>0$, dus er is een $N$ met $|u_n-L|<\dfrac{\varepsilon}{|k|}$ voor elke $n>N$. Voor die $n$ geldt $|ku_n-kL|=|k|\cdot|u_n-L|<\varepsilon$.
\end{fullsolution}
\end{oefening}

\begin{oefening}[3]
\thema{waar of niet waar}Waar of niet waar? Bewijs kort, of geef een tegenvoorbeeld.
\begin{enumerate}
  \item Elke begrensde rij is convergent.
  \item Elke stijgende rij is divergent.
  \item Elke convergente rij is begrensd.
  \item Als $u_n + v_n$ convergent is, dan zijn $u_n$ en $v_n$ allebei convergent.
  \item Als $\displaystyle\lim_{n\to+\infty} u_n = 0$, dan is $\displaystyle\lim_{n\to+\infty} u_n v_n = 0$.
\end{enumerate}
\begin{solution}
\textbf{a.}~niet waar \quad \textbf{b.}~niet waar \quad \textbf{c.}~waar \quad \textbf{d.}~niet waar \quad \textbf{e.}~niet waar
\end{solution}
\begin{fullsolution}
\textbf{a.}~Niet waar: $(-1)^n$.\\
\textbf{b.}~Niet waar: $1-\tfrac1n$.\\
\textbf{c.}~Waar: neem $\varepsilon=1$. Vanaf een bepaalde $N$ liggen alle termen tussen $L-1$ en $L+1$, en er blijven maar eindig veel andere termen over.\\
\textbf{d.}~Niet waar: $u_n=(-1)^n$ en $v_n=-(-1)^n$ geven som 0.\\
\textbf{e.}~Niet waar: $u_n=\tfrac1n$ en $v_n=n$ geven product 1.
\end{fullsolution}
\end{oefening}

\begin{oefening}[3]
\thema{reeks: telescoperen}Gegeven: $u_i = \dfrac{1}{i(i+1)}$.
\begin{enumerate}
  \item Toon aan dat $u_i = \dfrac1i - \dfrac{1}{i+1}$.
  \item Schrijf $s_n = u_1 + u_2 + \cdots + u_n$ uit en vereenvoudig.
  \item Bereken $\tfrac12 + \tfrac16 + \tfrac{1}{12} + \tfrac{1}{20} + \cdots$
\end{enumerate}
\begin{solution}
\textbf{b.}~$s_n=1-\dfrac{1}{n+1}$ \quad \textbf{c.}~1
\end{solution}
\begin{fullsolution}
\textbf{a.}~$\dfrac1i-\dfrac{1}{i+1}=\dfrac{i+1-i}{i(i+1)}=\dfrac{1}{i(i+1)}$.\\
\textbf{b.}~$s_n=\left(1-\tfrac12\right)+\left(\tfrac12-\tfrac13\right)+\cdots+\left(\tfrac1n-\tfrac{1}{n+1}\right)=1-\dfrac{1}{n+1}$: bijna alles valt weg.\\
\textbf{c.}~De rij is $u_i$, en de som is de limiet van $s_n$: 1. Deze reeks is niet meetkundig, en toch heeft ze een som.
\end{fullsolution}
\end{oefening}
\end{oefeningen}

%----------------------------------------------------------------------------------------
% KOPPEL: verbanden met andere leerstof
%----------------------------------------------------------------------------------------
\opdrachtsectie{Koppel}

\lkr{Exponentiële functies, complexe getallen en matrices zijn gezien in het 5de jaar. Kijk na of het getal $e$ al ingevoerd is. Logica en kansrekenen staan ook in het 6de jaar: kijk na of ze al aan bod kwamen, of laat die oefeningen weg.}

\begin{oefeningen}
\begin{oefening}[3]
\koppel{discreet}{Logica en bewijzen} De definitie van de limiet heeft de vorm $\forall\, \varepsilon > 0,\ \exists\, N \in \mathbb{N},\ \forall\, n > N: \ |u_n - L| < \varepsilon$.
\begin{enumerate}
  \item Schrijf de ontkenning met kwantoren. Wat betekent ze in woorden?
  \item Bewijs daarmee dat 1 niet de limiet is van $u_n = (-1)^n$.
\end{enumerate}
\begin{solution}
\textbf{a.}~$\exists\,\varepsilon>0,\ \forall\,N\in\mathbb{N},\ \exists\,n>N:\ |u_n-L|\geqslant\varepsilon$
\end{solution}
\begin{fullsolution}
\textbf{a.}~$\exists\,\varepsilon>0,\ \forall\,N\in\mathbb{N},\ \exists\,n>N:\ |u_n-L|\geqslant\varepsilon$. In woorden: er is een band waar steeds opnieuw termen buiten vallen, hoe ver je ook gaat.\\
\textbf{b.}~Neem $\varepsilon=1$. Voor elke $N$ is er een oneven $n>N$, en dan is $|u_n-1|=|-1-1|=2\geqslant 1$.
\end{fullsolution}
\end{oefening}

\begin{oefening}[2]
\koppel{analyse}{Exponentiële functies} Je belegt \eur 1 aan 100\% per jaar. Je krijgt de intrest niet één keer, maar $n$ keer per jaar, telkens $\tfrac{100}{n}$\%. Na een jaar heb je $u_n = \left(1 + \tfrac1n\right)^n$.
\begin{enumerate}
  \item Bereken $u_1$, $u_{12}$, $u_{365}$ en $u_{1\,000\,000}$ op vier decimalen.
  \item Word je zo onbeperkt rijk? Welk getal herken je?
\end{enumerate}
\begin{solution}
\textbf{a.}~2; 2,6130; 2,7146; 2,7183 \quad \textbf{b.}~nee: de limiet is $\mathrm{e}$
\end{solution}
\begin{fullsolution}
\textbf{a.}~$2^1=2$, $u_{12}\approx 2{,}6130$, $u_{365}\approx 2{,}7146$, $u_{1\,000\,000}\approx 2{,}7183$.\\
\textbf{b.}~Nee: de waarden stijgen, maar komen nooit boven $\mathrm{e}\approx 2{,}71828$. De rij is stijgend en naar boven begrensd.
\end{fullsolution}
\end{oefening}

\begin{oefening}[2]
\koppel{algebra}{Complexe getallen} Gegeven: $z = \tfrac12 + \tfrac12\mathrm{i}$.
\begin{enumerate}
  \item Bereken $|z|$. Wat is de limiet van $|z^n|$? Wat betekent dat voor $z^n$ in het complexe vlak?
  \item Bereken $1 + z + z^2 + \cdots$ met de formule van de meetkundige reeks.
\end{enumerate}
\begin{solution}
\textbf{a.}~$\tfrac{\sqrt2}{2}$; 0; de punten spiraleren naar de oorsprong \quad \textbf{b.}~$1+\mathrm{i}$
\end{solution}
\begin{fullsolution}
\textbf{a.}~$|z|=\sqrt{\tfrac14+\tfrac14}=\tfrac{\sqrt2}{2}<1$, dus $|z^n|=|z|^n\to 0$. De punten $z^n$ draaien telkens 45° en komen steeds dichter bij de oorsprong.\\
\textbf{b.}~$\dfrac{1}{1-z}=\dfrac{1}{\tfrac12-\tfrac12\mathrm{i}}=\dfrac{2}{1-\mathrm{i}}=\dfrac{2(1+\mathrm{i})}{2}=1+\mathrm{i}$.
\end{fullsolution}
\end{oefening}

\begin{oefening}[2]
\begin{metprent}{bakker}
\koppel{algebra}{Matrices} In een dorp zijn er twee bakkers. Elke maand stapt 10\% van de klanten van A over naar B, en 20\% van de klanten van B naar A. De aantallen klanten zijn $X_n = \begin{pmatrix} a_n \\ b_n \end{pmatrix}$, met $X_{n} = M X_{n-1}$ en $M = \begin{pmatrix} 0{,}9 & 0{,}2 \\ 0{,}1 & 0{,}8 \end{pmatrix}$.
\begin{enumerate}
  \item Bereken $X_1$ en $X_2$ als $X_0 = \begin{pmatrix} 300 \\ 0 \end{pmatrix}$.
  \item Zoek het evenwicht: los $X = MX$ op met $a + b = 300$. Vergelijk met $L = aL + b$.
\end{enumerate}
\end{metprent}
\begin{solution}
\textbf{a.}~$\begin{pmatrix} 270 \\ 30 \end{pmatrix}$, $\begin{pmatrix} 249 \\ 51 \end{pmatrix}$ \quad \textbf{b.}~200 en 100
\end{solution}
\begin{fullsolution}
\textbf{a.}~$X_1=\begin{pmatrix} 270 \\ 30 \end{pmatrix}$ en $X_2=\begin{pmatrix} 243+6 \\ 27+24 \end{pmatrix}=\begin{pmatrix} 249 \\ 51 \end{pmatrix}$.\\
\textbf{b.}~$a=0{,}9\,a+0{,}2\,b$ geeft $0{,}1\,a=0{,}2\,b$, dus $a=2b$. Met $a+b=300$ is $a=200$ en $b=100$. Het is dezelfde redenering, met een matrix in plaats van een getal: de limiet verandert niet meer bij nog een stap.
\end{fullsolution}
\end{oefening}

\begin{oefening}[2]
\koppel{data}{Kansrekenen} Twee spelers gooien om beurt met een dobbelsteen. Wie eerst een zes gooit, wint. Speler A begint.
\begin{enumerate}
  \item Bereken de kans dat A wint in zijn eerste, tweede en derde beurt. Welke soort rij is dat?
  \item Bereken de kans dat A het spel wint.
  \item Is het voordeel van beginnen groot?
\end{enumerate}
\begin{solution}
\textbf{a.}~$\tfrac16$, $\tfrac{25}{216}$, $\tfrac{625}{7776}$; meetkundig, $q=\tfrac{25}{36}$ \quad \textbf{b.}~$\tfrac{6}{11}$ \quad \textbf{c.}~ongeveer 54,5\% tegen 45,5\%
\end{solution}
\begin{fullsolution}
\textbf{a.}~Voor elke volgende beurt moeten beide spelers eerst geen zes gooien: telkens maal $(\tfrac56)^2=\tfrac{25}{36}$.\\
\textbf{b.}~$\dfrac{\tfrac16}{1-\tfrac{25}{36}}=\dfrac16\cdot\dfrac{36}{11}=\dfrac{6}{11}$.\\
\textbf{c.}~A wint met kans $\tfrac{6}{11}\approx 54{,}5\%$ en B met $\tfrac{5}{11}\approx 45{,}5\%$.
\end{fullsolution}
\end{oefening}
\end{oefeningen}

%----------------------------------------------------------------------------------------
\opdrachtsectie{Test jezelf}

\begin{oefeningen}
\begin{oefening}[1]\label{oef:h4-t1}
Bepaal de limiet. Is de rij convergent of divergent?
\begin{reeks}
  \item $u_n = 3 - \tfrac4n$
  \item $u_n = (-1)^n \cdot n$
  \item $u_n = 2n + 1$
\end{reeks}
\begin{solution}
\textbf{a.}~3, conv. \quad \textbf{b.}~bestaat niet, div. \quad \textbf{c.}~$+\infty$, div.
\end{solution}
\begin{fullsolution}
\textbf{a.}~$\tfrac4n\to 0$.\\
\textbf{b.}~De termen wisselen van teken en worden steeds groter in absolute waarde.\\
\textbf{c.}~De termen stijgen onbegrensd.
\end{fullsolution}
\end{oefening}

\begin{oefening}[1]\label{oef:h4-t2}
Bepaal de limiet.
\begin{reeks}
  \item $5,\ 10,\ 20,\ \ldots$
  \item $16,\ {-8},\ 4,\ \ldots$
  \item $7,\ 4,\ 1,\ \ldots$
\end{reeks}
\begin{solution}
\textbf{a.}~$+\infty$ \quad \textbf{b.}~0 \quad \textbf{c.}~$-\infty$
\end{solution}
\begin{fullsolution}
\textbf{a.}~$q=2$ en $u_1>0$.\\
\textbf{b.}~$q=-\tfrac12$, dus $|q|<1$.\\
\textbf{c.}~$v=-3$.
\end{fullsolution}
\end{oefening}

\begin{oefening}[2]\label{oef:h4-t3}
Bewijs met de definitie: $\displaystyle\lim_{n\to+\infty} \dfrac{3n+2}{n} = 3$.
\begin{solution}
Kies $N\geqslant\tfrac{2}{\varepsilon}$.
\end{solution}
\begin{fullsolution}
Neem $\varepsilon>0$. $\left|\dfrac{3n+2}{n}-3\right|=\dfrac2n<\varepsilon$ geeft $n>\dfrac{2}{\varepsilon}$. Kies $N\geqslant\dfrac{2}{\varepsilon}$: dan geldt voor elke $n>N$ dat $|u_n-3|<\varepsilon$.
\end{fullsolution}
\end{oefening}

\begin{oefening}[1]\label{oef:h4-t4}
Bereken $\displaystyle\lim_{n\to+\infty} \dfrac{2n^2 - n}{5n^2 + 3}$.
\begin{solution}
$\tfrac25$
\end{solution}
\begin{fullsolution}
Deel teller en noemer door $n^2$: $\dfrac{2-\tfrac1n}{5+\tfrac{3}{n^2}}\to\dfrac25$.
\end{fullsolution}
\end{oefening}

\begin{oefening}[2]\label{oef:h4-t5}
Bereken $\displaystyle\lim_{n\to+\infty} \dfrac{2 + \sin n}{n^2}$.
\begin{solution}
0
\end{solution}
\begin{fullsolution}
$1\leqslant 2+\sin n\leqslant 3$, dus $\dfrac{1}{n^2}\leqslant u_n\leqslant\dfrac{3}{n^2}$. Beide grenzen hebben limiet 0.
\end{fullsolution}
\end{oefening}

\begin{oefening}[2]\label{oef:h4-t6}
Toon aan dat $u_n = \dfrac{3n}{n+1}$ stijgend en naar boven begrensd is. Wat is de limiet?
\begin{solution}
Bovengrens 3; limiet 3
\end{solution}
\begin{fullsolution}
$u_n=3-\dfrac{3}{n+1}<3$. $u_{n+1}-u_n=\dfrac{3}{n+1}-\dfrac{3}{n+2}=\dfrac{3}{(n+1)(n+2)}>0$. De limiet is 3.
\end{fullsolution}
\end{oefening}

\begin{oefening}[1]\label{oef:h4-t7}
Gegeven: $u_0 = 100$ en $u_n = 0{,}6\,u_{n-1} + 20$. Bepaal het evenwicht.
\begin{solution}
50
\end{solution}
\begin{fullsolution}
$L=0{,}6\,L+20$ geeft $0{,}4\,L=20$, dus $L=50$.
\end{fullsolution}
\end{oefening}

\begin{oefening}[1]\label{oef:h4-t8}
Bereken $20 + 5 + \tfrac54 + \cdots$ en schrijf $0{,}363636\ldots$ als een breuk.
\begin{solution}
$\tfrac{80}{3}$ en $\tfrac{4}{11}$
\end{solution}
\begin{fullsolution}
$q=\tfrac14$, dus $\dfrac{20}{1-\tfrac14}=\dfrac{80}{3}$. Verder is $\dfrac{0{,}36}{1-0{,}01}=\dfrac{36}{99}=\dfrac{4}{11}$.
\end{fullsolution}
\end{oefening}
\end{oefeningen}

\herhaalschema{%
  \fv{oef:h4-t1}{vb:h4-soort}{oef:h4-soort}%
  \fv{oef:h4-t2}{vb:h4-rrmr}{oef:h4-rrmr}%
  \fv{oef:h4-t3}{vb:h4-bewijs}{oef:h4-bewijs}%
  \fv{oef:h4-t4}{vb:h4-reken}{oef:h4-ratio}%
  \fv{oef:h4-t5}{vb:h4-insluit}{oef:h4-insluit}%
  \fv{oef:h4-t6}{vb:h4-monotoon}{oef:h4-monotoon}%
  \fv{oef:h4-t7}{vb:h4-dyn}{oef:h4-evenwicht}%
  \fv{oef:h4-t8}{vb:h4-reeks}{oef:h4-reeks}}

\ifzelfstudie
\section*{Maak je eigen samenvatting}
Sluit dit boek en schrijf in je schrift: wat convergent en divergent betekent, de formele definitie, wanneer een rekenkundige en een meetkundige rij convergeren, hoe je een limiet berekent, de insluitstelling, de eigenschap van stijgende en begrensde rijen, hoe je een evenwicht vindt, en de som van een meetkundige reeks. Vergelijk daarna met de kaders.
\fi
